\documentclass[10pt,twocolumn]{article}
\usepackage[T1]{fontenc}
\usepackage[utf8]{inputenc}
\usepackage{lmodern}
\usepackage[margin=1.8cm,top=2.4cm,bottom=2cm,columnsep=0.6cm]{geometry}
\usepackage{fancyhdr}
\usepackage{titlesec}
\usepackage{booktabs}
\usepackage{amsmath,amssymb,amsfonts}
\usepackage{microtype}
\usepackage{xcolor}
\usepackage{mdframed}
\usepackage{parskip}
\usepackage{enumitem}
\usepackage{hyperref}

\definecolor{rulered}{RGB}{180,30,30}
\definecolor{boxgray}{RGB}{245,245,245}
\definecolor{keywordgray}{RGB}{100,100,100}

\pagestyle{fancy}
\fancyhf{}
\fancyhead[L]{\textsc{\small Claude Mag}}
\fancyhead[C]{\small\textit{Machine Learning Research Digest}}
\fancyhead[R]{\small 2026-08-31}
\fancyfoot[C]{\small\thepage}
\renewcommand{\headrulewidth}{0.8pt}

\titleformat{\section}{\normalsize\bfseries\color{rulered}}{}{0pt}{}%
  [\vspace{-4pt}\color{rulered}\rule{\columnwidth}{0.4pt}]
\titlespacing{\section}{0pt}{8pt}{4pt}

\hypersetup{colorlinks=true,urlcolor=rulered,linkcolor=black}

\begin{document}
\setlength{\parindent}{0pt}
\setlength{\parskip}{4pt}

{\small\color{keywordgray}\textbf{Keywords:}
  evolution strategies \textbullet{}
  LLM post-training \textbullet{}
  pass@k \textbullet{}
  solution diversity}\par\vspace{4pt}

{\LARGE\bfseries\raggedright
  Understanding Evolution Strategies for LLM
  Reasoning: Broader Coverage than GRPO}\par\vspace{4pt}

{\large\itshape\raggedright
  Population-Based Weight-Space Optimization Preserves Output
  Diversity That RL Post-Training Collapses}\par\vspace{6pt}

{\small\textbf{By} Yunpeng Ba, Zhi Zheng, Yue Xie, Jiaqing Li,
  Xialiang Tong, Tao Zhong, Mingxuan Yuan, Zhichao Lu,
  Xuyang Wu, Zhenkun Wang}\par\vspace{2pt}

{\small\color{keywordgray}arXiv:2608.27351 \textbullet{} Preprint, August 2026}
\par\vspace{2pt}\noindent{\color{rulered}\rule{\columnwidth}{0.6pt}}\vspace{6pt}

Post-training LLMs with reinforcement learning — most commonly via
GRPO — improves single-sample accuracy (pass@1) but systematically
narrows the model's output distribution, reducing solution diversity
across multiple samples. This paper formalises why Evolution Strategies
(ES), a gradient-free optimizer that perturbs \emph{weights} rather than
policy outputs, avoids this collapse: ES maintains the distributional
breadth that RL destroys. On reasoning benchmarks, ES matches GRPO on
pass@1 while outperforming it by up to $+14.2$ pp on pass@32.

\begin{mdframed}[backgroundcolor=boxgray,linecolor=rulered,linewidth=1pt,
    innerleftmargin=6pt,innerrightmargin=6pt,
    innertopmargin=6pt,innerbottommargin=6pt]
\textbf{\small AT A GLANCE}\par\vspace{4pt}
\begin{description}[leftmargin=0pt,labelsep=4pt,itemsep=2pt,parsep=0pt,
    font=\small\bfseries]
  \item[Problem:] RL post-training (GRPO) improves pass@1 by concentrating
    output probability mass, but this collapses pass@k diversity — the
    opposite of what discovery-oriented applications need.
  \item[Why it matters:] In math, science, and code, generating $k$
    diverse candidate solutions and checking the best is often more
    effective than optimising a single best guess.
  \item[Method:] Evolution Strategies perturb model weights via isotropic
    Gaussian noise, estimate a smoothed reward gradient over a population,
    and update parameters without ever conditioning on per-prompt output
    distributions — so entropy cannot collapse.
  \item[Result:] ES matches GRPO within 1--2 pp on pass@1; outperforms
    by $+6.3$ pp (MATH-500) to $+14.2$ pp (AMC/AIME) on pass@32.
  \item[Evidence:] ES-trained models retain 94\% of base model output
    entropy vs.\ 71\% for GRPO; unique solution fingerprints are 2.3$\times$
    more frequent at $k{=}8$.
\end{description}
\end{mdframed}

\section{The Big Picture}

Post-training is now the dominant mechanism for improving LLM reasoning.
Group Relative Policy Optimisation (GRPO) and its relatives reward
high-scoring outputs and penalise low-scoring ones, using group
statistics as a variance-reduction baseline. The resulting policy
update increases the log-probability of high-advantage outputs — and
in doing so, mechanically compresses the output distribution.

For pass@1 tasks — pick the single best answer — this is fine.
For pass@k tasks — check at least one of $k$ samples — compressed
distributions are harmful: the model generates variations of the
same high-reward solution rather than exploring the solution space.

Evolution Strategies sidestep this entirely by never touching output
distributions. Instead, ES evaluates perturbed \emph{weight} vectors
and updates the base parameters in the direction of higher-reward
perturbations. The output distribution at any fixed $\theta$ is
unchanged by the ES mechanism; only the weight location shifts.

\section{Why Existing Approaches Fall Short}

\textbf{GRPO} samples $G$ rollouts per prompt, computes group-relative
advantages $\hat{A}_i = (R_i - \bar{R})/\sigma_R$, and applies:
\[
  \nabla_\theta \mathcal{L}_\text{GRPO} = -\sum_i \hat{A}_i
    \nabla_\theta \log \pi_\theta(a_i | x)
\]
This directly increases $\log \pi_\theta(a_i|x)$ for high-reward
outputs, reducing $H(\pi_\theta(\cdot|x))$ each step. Over training,
output entropy converges to a low plateau regardless of the base model's
initial diversity. Ablations confirm this: GRPO trained for $T$ steps
reduces entropy to 71\% of the base model's, independent of $T$
beyond early training.

\textbf{PPO} has the same problem: the clipped probability-ratio
objective $\min(\rho_t \hat{A}_t,\, \text{clip}(\cdot)\hat{A}_t)$
increases log-probabilities of advantaged actions, compressing
distribution mass similarly.

\textbf{SFT} also concentrates distributions, but only on
demonstration patterns — it does not optimise diversity directly.

\section{Core Mechanism}

ES treats the model as a black box and optimises in weight space.
At iteration $t$, with current parameters $\theta_t$:

\begin{enumerate}[itemsep=2pt,parsep=0pt]
  \item Sample $N$ noise vectors $\epsilon_i \sim \mathcal{N}(0, I)$
  \item Evaluate perturbed models: $R_i = R(\theta_t + \sigma \epsilon_i)$
  \item Update:
    $\theta_{t+1} = \theta_t + \frac{\alpha}{N\sigma}\sum_i R_i \epsilon_i$
\end{enumerate}

The update is an unbiased estimate of
$\nabla_\theta \mathbb{E}_\epsilon[R(\theta + \sigma\epsilon)]$,
the gradient of the \emph{smoothed} objective.

Crucially, neither the update rule nor the reward evaluation
conditions on any per-prompt output distribution $\pi_\theta(\cdot|x)$.
Output entropy is regulated by the landscape of weight perturbations,
not forced to collapse by logit reinforcement.

\section{Technical Formulation}

The ES gradient estimate over population size $N$ is:
\[
  \hat{g}_\text{ES} = \frac{1}{N\sigma}\sum_{i=1}^N R_i \epsilon_i
\]
with variance $\text{Var}(\hat{g}) = \frac{\mathbb{E}[R^2]}{N\sigma^2}$,
decreasing in $N$. The paper shows that the smoothed objective
$F_\sigma(\theta) = \mathbb{E}_\epsilon[R(\theta + \sigma\epsilon)]$
satisfies:
\[
  F_\sigma(\theta) \geq F_0(\theta) - \frac{\sigma^2}{2}L_R
\]
where $L_R$ is the Lipschitz constant of $R$, bounding the gap
between the smoothed and original objectives.

For a fixed prompt $x$, the output entropy $H(\pi_{\theta+\sigma\epsilon}(\cdot|x))$
averages over the perturbation distribution rather than being
targeted directly. The expected entropy after one update is:
\[
  \mathbb{E}[H(\pi_{\theta+\hat{g}}(\cdot|x))] \approx H(\pi_\theta(\cdot|x))
    + O(\alpha \sigma \|\nabla H\|)
\]
compared to the GRPO decrease of $O(\alpha \sum_i \hat{A}_i)$
per step, which is systematically negative for positive-reward outputs.

\section{Learning or Inference Procedure}

Training with ES proceeds as:
\begin{enumerate}[itemsep=2pt,parsep=0pt]
  \item For each iteration, sample a batch of reasoning problems
  \item Sample $N=32$ noise vectors, evaluate all perturbed models in
    parallel on the same batch
  \item Compute ES gradient and update $\theta$
  \item Anneal $\sigma$ by a factor of 0.995 per iteration
\end{enumerate}

No rollout groups, no critics, no probability-ratio clipping.
At inference, only $\theta$ is used; the noise ensemble is discarded.

\section{What the Guarantee Says}

The paper proves that ES converges to a stationary point of the
smoothed objective $F_\sigma$ at rate $O(1/\sqrt{T})$ under standard
smoothness assumptions. The output-entropy preservation result is
proven for a single gradient step and validated empirically over
full training runs.

\section{Experimental Findings}

Evaluated with Qwen2.5-7B on MATH-500, AMC/AIME, and LiveCodeBench:

\begin{itemize}[itemsep=2pt,parsep=0pt]
  \item \textbf{pass@1:} ES within 1--2 pp of GRPO across all benchmarks
  \item \textbf{pass@8:} ES outperforms GRPO by $+6.3$ pp (MATH-500),
    $+11.4$ pp (AMC/AIME), $+8.7$ pp (LiveCodeBench)
  \item \textbf{pass@32:} Advantages grow further; AMC/AIME gap reaches
    $+14.2$ pp
  \item \textbf{Output entropy:} ES retains 94\% of base model entropy;
    GRPO retains only 71\%
  \item \textbf{Solution diversity:} 2.3$\times$ more unique solution
    fingerprints under ES at $k{=}8$
\end{itemize}

\section{Ablations and Interpretation}

\textbf{Population size $N$:} Larger $N$ reduces gradient variance
and improves both pass@1 and pass@k, with diminishing returns beyond
$N{=}32$. Smaller $N$ ($<$16) produces unstable training.

\textbf{Perturbation scale $\sigma$:} Larger $\sigma$ increases
exploration (better pass@k) at the cost of slower convergence on
pass@1. $\sigma{=}0.05$ provides a good balance.

\textbf{Comparison to SFT:} SFT with demonstration data improves
pass@1 comparably to GRPO but also degrades pass@k, confirming
that distribution narrowing is a general property of supervised
objectives, not just RL.

\textbf{Compute cost:} ES requires $N$ forward passes per iteration
vs.\ $G$ for GRPO. At $N{=}G{=}32$, wall-clock time is comparable
because GRPO requires both forward and backward passes while ES
requires only forward passes (no gradient through the model).

\section*{Reference}

Yunpeng Ba, Zhi Zheng, Yue Xie, Jiaqing Li, Xialiang Tong, Tao Zhong,
Mingxuan Yuan, Zhichao Lu, Xuyang Wu, Zhenkun Wang.
\textbf{Understanding Evolution Strategies for LLM Reasoning:
Broader Reasoning Coverage than GRPO.}
arXiv:2608.27351, August 2026.
\url{https://arxiv.org/abs/2608.27351}

\end{document}
